\section{Planned Extensions}
\label{sec:planned-extensions}

This section tracks methodology changes that have been decided on but not yet
executed, identified while reviewing this report's description of the pipeline against
what the pipeline should actually do. It is expected to grow as further review surfaces
further such directions; each entry stays here until the corresponding pipeline change
has been made and the rest of this report updated to reflect it.

\subsection{Unify Hardware-pillar corpus construction with the common procedure}
\label{sec:planned-hardware-unification}

As noted in Section~\ref{sec:candidate-generation}, the quantum-side Hardware pillar is
currently built by a separate, older two-source procedure (arXiv plus Crossref
cross-referencing against seven target journals) rather than the common per-pillar
arXiv query procedure used by every other pillar on both sides --- an inconsistency
left over from before the six-pillar comparison architecture existed, not a deliberate
design choice. The two-source method is, if anything, more thorough than the common
procedure (it additionally catches journal-only publications with no arXiv preprint),
so the fix is to extend it outward rather than to drop it: the Crossref
cross-referencing step will be added to the common per-pillar procedure
(Section~\ref{sec:pillar-corpus-construction}) for every pillar on both sides, using
each pillar's own keyword list and an appropriately chosen set of target journals,
rather than rebuilding the Hardware pillar under the weaker arXiv-only method it was
originally exempted from.

This extension will be applied even to pillars, and specifically to the quantum
Hardware pillar, that already have completed classification work
(Table~\ref{tab:corpus-stats}). Any additional papers the Crossref cross-referencing
step finds beyond what is already in a pillar's corpus will be deduplicated against the
existing corpus by arXiv identifier (and, for journal-only records with no arXiv
identifier, by the same fuzzy title-matching procedure already used in
Section~\ref{sec:quantum-list-gen}) and added to it; no paper already in a corpus, and
no classification already completed against it, will be removed or discarded as a
result of this extension. This follows the same additive-only principle already
applied to the keyword-list rework (Section~\ref{sec:keyword-rework}): a corpus-widening
change should only ever add coverage, never retract it.

\subsection{Decide whether CS needs an out-of-scope category analogous to quantum's}
\label{sec:planned-cs-out-of-scope}

Section~\ref{sec:domain-classification} resolves the quantum side's
computing-vs-physics scope question by classifying out-of-scope papers into an
existing set of non-computing labels (\emph{Quantum Sensing}, \emph{Atomic Physics},
etc.) at the domain-classification stage, rather than restricting the corpus-pull
query and losing recall. Whether the CS side has an analogous problem --- a keyword
matching a paper that is not really about the intended pillar's subject at all, as
opposed to a paper that is on-topic but assigned to the wrong one of the six pillars
--- has not yet been assessed. A candidate example: ``channel capacity'' and
``information theoretic limit'' (CS info\_theory) are both used extensively in pure
information-theory mathematics and in communications-engineering work with no
connection to computing. This needs the same kind of review the quantum side already
received before deciding whether CS's domain classifier needs its own out-of-scope
label set, and if so, what it should contain.
